
\section{Congruence classes and elementary collapses}

\begin{lemma}\label{ptc-kernel:lem:lower}
If $w$ is a word of length $\ell$, then
\[
 xw\equiv x-\ell\pmod d\qquad(x\in\mathcal N).
\]
Moreover $Cw\subseteq C$ and $|Cw|\ge d$. In particular, $\rank(w)\ge d$.
\end{lemma}
\begin{proof}
For every $i\in I$, we have $xM_i\equiv x-1\pmod d$: at $x=1$ this
uses $d\mid i$, and at all other points it follows from the definition.
Induction proves the formula for words. Since every generator preserves
$C$, so does every word. The set $C$ meets every residue class modulo
$d$, and the congruence formula shows that its image does as well.
\end{proof}

\begin{lemma}\label{ptc-kernel:lem:collapse}
For each $i\in I$ with $i<m$, and each $a\in C$, the restriction
semigroup $S|_C$ contains the transformation fixing every point other
than $a$ and sending $a$ to $a+i$ modulo $m$.
\end{lemma}
\begin{proof}
On $C$, let $c=M_m|_C$. This is the permutation $x\mapsto x-1$
modulo $m$, so $c^{-1}=c^{m-1}$. The restriction of
$M_iM_m^{m-1}$ fixes every $x>1$ and sends $1$ to $i+1$.
Conjugation by powers of $c$ gives all the asserted transformations.
These conjugations are realized by words in the original generators,
since both $c$ and $c^{-1}$ are restrictions of positive powers of $M_m$.
\end{proof}

\begin{lemma}\label{ptc-kernel:lem:retraction}
For every $A\in\mathcal A$, there exists $v_A\in S$ such that
\[
 xv_A=a_r\quad(x\in C_r),\qquad
 xv_A\equiv x\pmod d\quad(x\in\mathcal N).
\]
\end{lemma}
\begin{proof}
Form the directed graph on $\mathbb{Z}/m\mathbb{Z}$ with edges
\[
 a\longrightarrow a+i\qquad(i\in I,\ i<m).
\]
Its strongly connected components are exactly the classes modulo $d$.
Indeed, the increments generate the subgroup $d\mathbb{Z}/m\mathbb{Z}$. In a finite
cyclic group the additive inverse of any increment is a nonnegative
multiple of that increment, so directed reachability gives precisely
the cosets of this subgroup.

In each component choose a directed spanning tree oriented toward the
prescribed root $a_r$. Such a tree exists, for example by choosing from
each nonroot vertex an edge along a shortest directed path to the root.
For each nonroot vertex apply the elementary collapse to its parent,
in decreasing order of distance from the root. A point first moves to
its parent and is subsequently moved along the remaining path to the
root. Thus the product of these collapses is the desired retraction
on $C$.

Choose words realizing these elementary collapses as in
Lemma~\ref{ptc-kernel:lem:collapse}, and concatenate them to obtain $v_A$.
Each collapse preserves residues modulo $d$. By
Lemma~\ref{ptc-kernel:lem:lower}, a representing word for such a collapse must
have length divisible by $d$ (apply the congruence formula to any point
of $C$). Hence that word preserves residues on all of $\mathcal N$,
not only on $C$. The product $v_A$ does too.

If there are no nonroot vertices, then $d=m$, $I=\{m\}$, and $A=C$.
In this case take $v_A=M_m^m$, which has the required properties.
\end{proof}
